\documentclass[aspectratio=169,12pt]{beamer}
\usepackage{fontspec}
\usepackage[portuguese]{babel}
\usepackage{xcolor,listings,tikz,booktabs,amsmath,amssymb,graphicx,multicol,array,colortbl}
\usetikzlibrary{arrows.meta,positioning,shapes.geometric,calc}

\definecolor{navy}{HTML}{0D1B3E}
\definecolor{gold}{HTML}{F5A623}
\definecolor{qgreen}{HTML}{1E8B4C}
\definecolor{qblue}{HTML}{1A6EAE}
\definecolor{codebg}{HTML}{EEF3F8}

\usetheme{default}
\useinnertheme{circles}
\setbeamertemplate{navigation symbols}{}
\setbeamertemplate{footline}{%
  \hfill{\color{gray!60}\scriptsize\insertframenumber/\inserttotalframenumber}\hspace{.6em}\vspace{.4em}}

\setbeamercolor{frametitle}{bg=navy,fg=white}
\setbeamercolor{structure}{fg=navy}
\setbeamercolor{alerted text}{fg=gold}
\setbeamercolor{item}{fg=gold}
\setbeamercolor{subitem}{fg=navy}
\setbeamercolor{block title}{bg=qblue,fg=white}
\setbeamercolor{block body}{bg=qblue!8,fg=black}
\setbeamercolor{block title alerted}{bg=gold,fg=navy}
\setbeamercolor{block body alerted}{bg=gold!15,fg=navy}
\setbeamercolor{block title example}{bg=qgreen,fg=white}
\setbeamercolor{block body example}{bg=qgreen!10,fg=black}
\setbeamerfont{frametitle}{size=\large,series=\bfseries}
\setbeamerfont{block title}{series=\bfseries}

\setbeamertemplate{title page}{%
  \begin{tikzpicture}[remember picture,overlay]
    \fill[navy](current page.north west)rectangle(current page.south east);
    \fill[gold]([yshift=.7cm]current page.south west)rectangle(current page.south east);
    \node[anchor=south,yshift=.73cm,navy,font=\tiny\bfseries]at(current page.south)
      {INTENSIVÃO QUANT AI \textbullet{} IMPACT UFSCAR};
  \end{tikzpicture}
  \vspace{.65cm}
  \begin{center}
    {\color{gold}\scriptsize\bfseries INTENSIVÃO QUANT AI \textbullet{} IMPACT UFSCAR}\\[.4cm]
    {\color{white}\LARGE\bfseries\inserttitle}\\[.25cm]
    {\color{gold!85}\normalsize\insertsubtitle}\\[.5cm]
    {\color{white!70}\small\insertauthor}\\[.1cm]
    {\color{white!50}\footnotesize\insertdate}
  \end{center}}

\lstdefinestyle{python}{language=Python,
  basicstyle=\ttfamily\scriptsize\color{qblue},
  keywordstyle=\color{navy}\bfseries,
  stringstyle=\color{qgreen!80!black},
  commentstyle=\color{gray!70}\itshape,
  backgroundcolor=\color{codebg},
  frame=l,framerule=2pt,rulecolor=\color{gold},
  xleftmargin=1em,breaklines=true,showstringspaces=false,tabsize=4,
  extendedchars=false,inputencoding=ascii}
\lstset{style=python}

\author{Felipe Maldonado\\{\scriptsize VP Diretoria Quant~\textbullet{}~Impact UFSCAR}}
\date{Intensivão Quant AI 2025}
\title{Aula 09 --- GenAI}
\subtitle{LLMs e a API Claude no Pipeline Quantitativo}

\begin{document}

\begin{frame}[plain]
\titlepage
\end{frame}

\begin{frame}[fragile]{Nesta Aula}
\begin{columns}[T]
  \column{0.5\textwidth}
\begin{block}{Teoria (25 min)}
\begin{itemize}
  \item LLMs no contexto de finanças quantitativas
  \item Anatomia de uma API de LLM: messages, tokens, temperature
  \item Prompt engineering para análise financeira
  \item Casos de uso: sentiment, narrativa, assistência a código
\end{itemize}
\end{block}
  \column{0.47\textwidth}
\begin{block}{Código (75 min)}
\begin{itemize}
  \item Setup: \texttt{anthropic} SDK, API key segura
  \item \texttt{chamar\_claude(prompt, system)} — wrapper genérico
  \item \texttt{gerar\_comentario\_performance()} — narrativa institucional
  \item \texttt{sugerir\_melhorias()} — pesquisa assistida por LLM
  \item Salvar narrativa em .txt para o relatório
\end{itemize}
\end{block}
\end{columns}
\end{frame}

\begin{frame}[fragile]{LLMs no Contexto Quant}
\begin{columns}[T]
  \column{0.5\textwidth}
\begin{itemize}
  \item \textbf{Análise de Sentimento}: processar earnings calls, atas do COPOM, releases — extrair score de sentimento com LLM
  \item \textbf{Geração de Narrativa}: dado Sharpe, MDD, IC $\Rightarrow$ comentário de performance no estilo gestora institucional
  \item \textbf{Assistência a Código}: descrever sinal em linguagem natural, LLM propõe implementação em Python
  \item \textbf{Pesquisa Assistida}: perguntar sobre literatura acadêmica, LLM cita papers e mecanismos relevantes
\end{itemize}
  \column{0.47\textwidth}
\begin{block}{Modelos Anthropic}
\begin{itemize}
  \item \textbf{Claude Haiku}: rápido, barato — experimentação
  \item \textbf{Claude Sonnet}: balanceado — produção
  \item \textbf{Claude Opus}: máxima capacidade — análises complexas
\end{itemize}
\end{block}\vspace{.3em}
\begin{alertblock}{Limitações importantes}
\begin{itemize}
  \item LLM não acessa internet nem dados em tempo real
  \item Alucinações: sempre verifique citações e números
  \item Contexto de treinamento: conhecimento até data de corte
  \item Custo por token: dimensione prompts com cuidado
\end{itemize}
\end{alertblock}
\end{columns}
\end{frame}

\begin{frame}[fragile]{Prompt Engineering Financeiro}
\begin{columns}[T]
  \column{0.5\textwidth}
\textbf{Prompt ruim:}
\begin{lstlisting}
"Analise minha estrategia."
\end{lstlisting}
\vspace{.2em}
\textbf{Prompt bom (estrutura completa):}
\begin{lstlisting}
system = '''Voce e um analista quantitativo
senior de uma gestora brasileira.
Estilo: tecnico, formal, preciso.
Tamanho: 200-250 palavras.'''

user = f'''Estrategia: momentum cross-sectional IBOVESPA.
Sharpe OOS: {sharpe:.2f}
CAGR: {cagr:.1%}
Max Drawdown: {mdd:.1%}
Escreva comentario de performance.'''
\end{lstlisting}
  \column{0.47\textwidth}
\begin{block}{Boas práticas}
\begin{itemize}
  \item System prompt: define personalidade e formato
  \item User prompt: dados concretos + instrução específica
  \item Temperature $\approx 0$: respostas determinísticas para números
  \item Temperature $\approx 0{,}3$: mais criatividade para narrativa
  \item few-shot: inclua 1--2 exemplos do formato desejado
\end{itemize}
\end{block}\vspace{.3em}
\begin{alertblock}{Segurança}
Nunca inclua API key no código. Use variável de ambiente: \texttt{ANTHROPIC\_API\_KEY}. Nunca suba \texttt{.env} para o GitHub.
\end{alertblock}
\end{columns}
\end{frame}

\begin{frame}[fragile]{O que Vamos Construir}
\begin{columns}[T]
  \column{0.52\textwidth}
\begin{block}{Funções a implementar}
\begin{itemize}
  \item \texttt{chamar\_claude(prompt, system, modelo, temp)} $\to$ str
  \item \texttt{resumo\_metricas(ret, nome)} $\to$ dicionário formatado
  \item \texttt{gerar\_comentario\_performance(m\_is, m\_oos)} $\to$ texto
  \item \texttt{sugerir\_melhorias(m\_oos)} $\to$ sugestões da literatura
\end{itemize}
\end{block}
  \column{0.45\textwidth}
\begin{block}{Entradas (parquets)}
\begin{itemize}
  \item \texttt{retorno\_carteira.parquet}
  \item \texttt{retorno\_walkforward\_liquido.parquet}
  \item API key da Anthropic configurada
\end{itemize}
\end{block}
\vspace{.4em}
\begin{exampleblock}{Saídas (parquets)}
\begin{itemize}
  \item \texttt{narrativa\_performance.txt}
  \item Sugestões de melhoria baseadas em literatura
\end{itemize}
\end{exampleblock}
\end{columns}
\end{frame}

\begin{frame}[fragile]{Referências}
\begin{multicols}{2}
\tiny
\begin{itemize}
  \item \textbf{Brown, T. et al.} (2020). Language models are few-shot learners (GPT-3). \textit{NeurIPS}.
  \item \textbf{Anthropic} (2024). Claude model card. \textit{anthropic.com}.
  \item \textbf{Wei, J. et al.} (2022). Chain-of-thought prompting elicits reasoning. \textit{NeurIPS}.
  \item \textbf{Lopez-Lira, A. \& Tang, Y.} (2023). Can ChatGPT forecast stock price movements? \textit{SSRN}.
  \item \textbf{Jegadeesh, N. \& Titman, S.} (1993). Returns to buying winners. \textit{Journal of Finance}, 48(1), 65--91.
\end{itemize}
\end{multicols}
\end{frame}


\end{document}
